\section{The \texttt{Lens\_CRL\_RTM} McXtrace Component}
1D CRL stack based on RTM formalism

\subsection*{Identification}
\begin{itemize}
  \item \textbf{Author:} Erik Knudsen
  \item \textbf{Origin:} DTU Physics
  \item \textbf{Date:} Jan '21
\end{itemize}

\subsection*{Description}
A CRL stack component based on the formalism presented by Simons et.al. J. Synch. Rad.

\begin{verbatim}
2017, vol. 24.
\end{verbatim}

We model a 1D lens stack focusing in the y-direction. I.e. invariant along x.

Example: Lens\_CRL\_RTM( r=0.5e-3, N=10, fast=1, yheight=0, d=0.1e-3,xwidth=1e-4, zdepth=4e-4)

\subsection*{Input parameters}
Parameters in \textbf{boldface} are required; the others are optional.

\begin{longtable}{p{0.22\textwidth}p{0.12\textwidth}p{0.46\textwidth}p{0.14\textwidth}}
\toprule
\textbf{Name} & \textbf{Unit} & \textbf{Description} & \textbf{Default} \\
\midrule
\endhead
r & m & Radius of curavature at the lens apex. & 0.5e-3 \\
d & m & Thickness of a single lens & 0.1e-3 \\
material & str & Datafile containing f1 constants & "Be.txt" \\
N & m & Number of lenslets in the stack. & 1 \\
zdepth & m & Thickness of a single lenslet. & 2e-3 \\
yheight & m & Height of lens opening. If zero this is set by the lens thickness. & 1e-3 \\
xwidth & m & Width of the lenslets. & 1.2e-3 \\
fast & m & Use fast calculation - should be off for better display with mxdisplay & 1 \\
\bottomrule
\end{longtable}

\subsection*{Links}
\begin{itemize}
  \item Component source code found in file \texttt{Lens\_CRL\_RTM.comp}.
\end{itemize}
\IfFileExists{optics/Lens_CRL_RTM_static.tex}{\input{optics/Lens_CRL_RTM_static.tex}}{}